% ---- Opci uvod (ne numerise se) --------------------------------------------
\chapter*{Uvod}
\addcontentsline{toc}{chapter}{Uvod}

Ugradbeni sistemi danas \v{c}ine najve\'ci dio ra\v{c}unarskih ure\dj{}aja u
svakodnevnoj upotrebi, od ku\'cnih aparata do industrijske opreme. Veliki broj
takvih sistema ima zahtjeve za vremenskom predvidljivo\v{s}\'cu, pa se za
njihovo upravljanje koriste operativni sistemi za rad u realnom vremenu. Jedan
od naj\v{c}e\v{s}\'ce kori\v{s}tenih takvih sistema otvorenog koda je FreeRTOS,
koji je dostupan za veliki broj arhitektura procesora. Istovremeno se sve vi\v{s}e
koristi otvorena skupovna arhitektura RISC-V, koja omogu\'cava izradu vlastitih
procesorskih jezgara bez licencnih ograni\v{c}enja. Spajanje ova dva otvorena
pristupa, operativnog sistema otvorenog koda i otvorene procesorske arhitekture,
postaje zanimljiva i aktuelna tema, posebno na jeftinim razvojnim plo\v{c}ama sa
programabilnom logikom.

Razvojna plo\v{c}a Tang Nano 9K, zasnovana na programabilnom logi\v{c}kom polju
(\foreignlanguage{english}{\emph{Field Programmable Gate Array}}, FPGA), pru\v{z}a
pristupa\v{c}an na\v{c}in da se na njoj sintetizuje meki procesor i izvr\v{s}ava
stvarni program. Meki procesor PicoRV32 jednostavna je RISC-V jezgra pogodna za
takve eksperimente. Me\dj{}utim, resursi ovakvih plo\v{c}a su veoma ograni\v{c}eni,
prije svega blokovska memorija i koli\v{c}ina logi\v{c}kih elemenata, pa
uklapanje slo\v{z}enijih sistema postaje in\v{z}enjerski izazov. Upravo to
\v{c}ini portanje cijele aplikacije zajedno sa operativnim sistemom i video
izlazom netrivijalnim zadatkom.

Ovaj rad nastavlja se na postoje\'cu aplikaciju TetriSaraj, koja je izvorno
pisana kao monolitni program bez operativnog sistema. Postavljaju se sljede\'ca
pitanja na koja rad treba odgovoriti. Da li se aplikacija mo\v{z}e prevesti u
skup zasebnih zadataka pod upravom raspore\dj{}iva\v{c}a operativnog sistema
FreeRTOS i dokazati njen rad na stvarnom hardveru? Da li se takav sistem,
zajedno sa video izlazom preko digitalnog multimedijalnog su\v{c}elja
(\foreignlanguage{english}{\emph{High-Definition Multimedia Interface}}, HDMI),
mo\v{z}e uklopiti u stroge resurse plo\v{c}e Tang Nano 9K? Te kona\v{c}no, kolika
je cijena zamjene konteksta na stvarnom hardveru i kako se ona odnosi prema
ranije dobijenim vrijednostima iz simulatora? Pretpostavlja se da je sve to
mogu\'ce uz pa\v{z}ljivo upravljanje memorijom i resursima.

Rad je organizovan na sljede\'ci na\v{c}in. U prvom poglavlju daje se pregled
srodnih radova i smje\v{s}ta vlastiti rad u kontekst postoje\'ce literature.
Drugo poglavlje opisuje hardversku platformu i prilago\dj{}eni port operativnog
sistema FreeRTOS na meki procesor, te osnovni dokaz njegovog rada. Tre\'ce
poglavlje prikazuje postupak razdvajanja aplikacije na zadatke i dokaz njihovog
rada na hardveru bez video izlaza. \v{C}etvrto poglavlje obja\v{s}njava problem
zauze\'ca resursa pri uklapanju cijele igre sa video izlazom i rje\v{s}enja koja
su to omogu\'cila. Peto poglavlje opisuje mjerenje vremena zamjene konteksta na
hardveru i pore\dj{}enje sa simulatorom. \v{S}esto poglavlje izdvaja mogu\'ce
pravce daljnjeg istra\v{z}ivanja. Na kraju rada dat je zaklju\v{c}ak u kojem su
sistematizirani najva\v{z}niji rezultati i doprinosi rada. Nakon zaklju\v{c}ka
nalazi se popis slika i lista kori\v{s}tenih bibliografskih referenci.
